\section{Anexos}

\subsection{Drive de Laboratorio}
La carpeta de Google Drive correspondiente a este laboratorio puede encontrarse en \url{https://drive.google.com/drive/folders/1zlPwmZ6XGEYrSYWA3p0l1uxA-KCTYJkX?usp=drive_link}. Allí encontrará evidencias del trabajo realizado durante esta actividad.

\subsection{Fragmentos de código}


\paragraph{SPI Init Maestro}\label{anexo:MasterSPIIn}---\\
\begin{lstlisting}[language=C, caption={Inicializacion SPI Maestro}]
void spi_init(void) {
	DDRB |= (1 << CS) | (1 << MOSI) | (1 << SCK); //Pines de comunicación
	DDRB &= ~(1 << MISO); // Entrada en Master in slave out
	SPCR = (1 << SPE) | (1 << MSTR); // SPI enable y Modo maestro
}
\end{lstlisting}

\paragraph{SPI Init Esclavo}\label{anexo:SlaveSPIIn}---\\
\begin{lstlisting}[language=C, caption={Inicializacion SPI Esclavo}]
void spi_init_slave(){
		DDRB |= (1 << MISO); // MISO salida
		DDRB &= ~(1 << MOSI); // MOSI entrada
		DDRB &= ~(1 << SCK); // Reloj de entrada, lo dicta el maestro	
		DDRB &= ~(1 << CS); // direccion habilitada por el maestro
		
		SPCR |= (1 << SPE); // Habilita SPI en modo maestro y no es necesario configurarle el prescaler, porque SCK se adapta al del maestro
}
\end{lstlisting}

\paragraph{Transferencia via SPI}\label{anexo:Transfer_spi}---\\
\begin{lstlisting}[language=C, caption={Transferencia SPI Maestro}]
uint8_t spi_transfer(uint8_t data) {
	SPDR = data;
	while (!(SPSR & (1 << SPIF)));
	return SPDR;
}
\end{lstlisting}

\paragraph{Recepcion del esclavo}\label{anexo:Reception_spi}---\\
\begin{lstlisting}[language=C, caption={Recepcion SPI Esclavo}]
uint8_t SPI_slaveReceive()
{
	// Register Clean
	SPDR = 0xFF;

	// Wait for reception complete SPIF(SPI Interrupt Flag)
	while(!(SPSR & (1 << SPIF)));

	// return Data Register, when is read SPIF is clean
	return SPDR;
}
\end{lstlisting}

\paragraph{Transferencia foto resistor}\label{anexo:fotoresistor}---\\
\begin{lstlisting}[language=C, caption={Transferencia foto resistor}]
uint8_t SPI_slaveReceive()
uint16_t lectura_led = adc_read_AC1();
		if (lectura_led > 650){
			SS_LOW(); spi_transfer(0xB6); SS_HIGH();
			strcpy(str_led,"Prendido");
			} else {
			SS_LOW(); spi_transfer(0xB7); SS_HIGH();
			strcpy(str_led,"Apagado");
		}
\end{lstlisting}


\paragraph{Transferencia de angulo del potenciometro}\label{anexo:potenciometro}---\\
\begin{lstlisting}[language=C, caption={Transferencia potenciometro}]
uint16_t pot = adc_read_AC0();
		float grados = pot * 0.17; // el 0.17 sale de 180/1024 - grados/adc
		Add_to_string(str_grados, (uint16_t)grados);
		SS_LOW(); spi_transfer((uint8_t)grados); SS_HIGH();
\end{lstlisting}

\paragraph{Transferencia de distancia}\label{anexo:distancia}---\\
\begin{lstlisting}[language=C, caption={Transferencia distancia}]

		PORTD |= (1<<PIN_T);
		_delay_us(15);
		PORTD &= ~(1<<PIN_T);

		if (Distancia_cm < 10){
			strcpy(str_dist,"Rojo");
			SS_LOW(); spi_transfer(0xB9); SS_HIGH();
		}
		else if (Distancia_cm < 40){
			strcpy(str_dist,"Naranja");
			SS_LOW(); spi_transfer(0xBA); SS_HIGH();
		}
		else if (Distancia_cm < 90){
			strcpy(str_dist,"Amarillo");
			SS_LOW(); spi_transfer(0xBB); SS_HIGH();
		}
		else if (Distancia_cm < 150){
			strcpy(str_dist,"Verde");
			SS_LOW(); spi_transfer(0xBC); SS_HIGH();
		}
		else if (Distancia_cm < 300){
			strcpy(str_dist,"Celeste");
			SS_LOW(); spi_transfer(0xBD); SS_HIGH();
		}
		else {
			strcpy(str_dist,"Azul");
			SS_LOW(); spi_transfer(0xBF); SS_HIGH();
		}


\end{lstlisting}


\paragraph{Interrupcion por comparacion hcsr04}\label{anexo:hcsr04}---\\
\begin{lstlisting}[language=C, caption={hcsr04}]


ISR(TIMER1_CAPT_vect) {
	if (TCCR1B & (1 << ICES1)) {
		start = ICR1;
		TCCR1B &= ~(1 << ICES1);
		} else {
		endd = ICR1;
		width = endd - start;
		TCCR1B |= (1 << ICES1);
		Distancia_cm = width / 116.0;
	}
}

\end{lstlisting}

\paragraph{Lectura DHT11}\label{anexo:DHT11}---\\
\begin{lstlisting}[language=C, caption={DHT11}]


bool dht11_read2(uint8_t *t,uint8_t *h){
	uint8_t d[5]={0};
		
	// Señal para datos	
	DDRD|=(1<<DHT_PIN);
	PORTD&=~(1<<DHT_PIN); 
	_delay_ms(20);
	PORTD|=(1<<DHT_PIN);
	 _delay_us(40);
	 
	// Entrada para datos, mismo pin
	DDRD&=~(1<<DHT_PIN); 
	
	PORTD|=(1<<DHT_PIN);
	
	uint16_t to=0;
	
	// Espera a la señal
	while(PIND&(1<<DHT_PIN)){ //Primer alto
		if(++to>200) 
		return false;
		_delay_us(1);
		}
		to=0;
		
	 while(!(PIND&(1<<DHT_PIN))){ // Primer bajo
		  if(++to>200)return false; 
		  _delay_us(1);}
		to=0; 
	while(PIND&(1<<DHT_PIN)){ // Segundo alto
		 if(++to>200)return false;
		_delay_us(1);}
	// Toma datos
	for(uint8_t i=0;i<40;i++){
		to=0; 
		while(!(PIND&(1<<DHT_PIN))){ 
			if(++to>200)return false; 
			_delay_us(1);}
		uint16_t w=0; while(PIND&(1<<DHT_PIN)){ 
			if(++w>255)break;
			_delay_us(1);
			}
		// Para ir guardando el bit 1 a 1 en d[i/8](8/8 = 1, 16/8 = 2, etc)
		d[i/8]<<=1; 
		if(w>40) 
		d[i/8]|=1;
	}
	if((uint8_t)(d[0]+d[1]+d[2]+d[3])!=d[4]) 
	return false;
	*h=d[0]; *t=d[2]; // Guardar Temperatura y Guardar Humedad
	return true;
}


\end{lstlisting}


\paragraph{Transferencia DHT11}\label{anexo:DHT11T}---\\
\begin{lstlisting}[language=C, caption={DHT11T}]
		if(dht11_read2(&Temp,&Hum) && leer_dht){
			leer_dht=0;
			if(Temp>26){
				SS_LOW(); spi_transfer(0xF0); SS_HIGH();
				} else {
				SS_LOW(); spi_transfer(0xB5); SS_HIGH();
			}
		}

		Add_to_string(str_temp, Temp);
		Add_to_string(str_hum, Hum);

\end{lstlisting}



\paragraph{Recepcion Esclavo}\label{anexo:RE}---\\
\begin{lstlisting}[language=C, caption={Recepcion Eclavo}]
    while(1)
    {
      
	      uint8_t byte = SPI_slaveReceive();

	      switch (byte) {

		      // ------------------------
		      // Control de PORTD2
		      // ------------------------
		      case 0xB6:
		      PORTD |= (1 << PORTD2);
		      break;

		      case 0xB7:
		      PORTD &= ~(1 << PORTD2);
		      break;

		     

		      // ------------------------
		      // Colores RGB
		      // ------------------------
		      case 0xB9: // Rojo
		      PORTD |=  (1 << PD3);
		      PORTD &= ~(1 << PD4);
		      PORTD &= ~(1 << PD5);
		      break;

		      case 0xBA: // Naranja (R+G)
		      PORTD |=  (1 << PD3);
		      PORTD |=  (1 << PD4);
		      PORTD &= ~(1 << PD5);
		      break;

		      case 0xBB: // Amarillo
		      PORTD |=  (1 << PD3);
		      PORTD |=  (1 << PD4);
		      PORTD &= ~(1 << PD5);
		      break;

		      case 0xBC: // Verde
		      PORTD &= ~(1 << PD3);
		      PORTD |=  (1 << PD4);
		      PORTD &= ~(1 << PD5);
		      break;

		      case 0xBD: // Celeste (G+B)
		      PORTD &= ~(1 << PD3);
		      PORTD |=  (1 << PD4);
		      PORTD |=  (1 << PD5);
		      break;

		      case 0xBF: // Azul
		      PORTD &= ~(1 << PD3);
		      PORTD &= ~(1 << PD4);
		      PORTD |=  (1 << PD5);
		      break;

		      // ------------------------
		      // Control PWM (OCR0A)
		      // ------------------------
		      case 0xF0:
		      OCR0A = 250;
		      break;

		      case 0xB5:
		      OCR0A = 0;
		      break;
			  
			// ------------------------
			// Servo (0 a 180)
			// ------------------------
			   
			  default:
			  if (byte < 180) {
				  deg = byte;
				  ticks = TMIN + ((uint32_t)(TMAX - TMIN) * deg) / 180;
				  OCR1A = ticks;
			  }
			  break;
	      }
      }

\end{lstlisting}


\paragraph{I2c en codigo de spi (Para LCD)}\label{anexo:LCDspi}---\\
\begin{lstlisting}[language=C, caption={I2c en spi}]

void I2C_init(void) {
	TWSR = 0x00;
	TWBR = ((F_CPU / 100000UL) - 16) / 2;
	TWCR = (1<<TWEN);
}

void I2C_start(void) {
	TWCR = (1<<TWINT)|(1<<TWSTA)|(1<<TWEN);
	while(!(TWCR&(1<<TWINT)));
}

void I2C_stop(void) {
	TWCR = (1<<TWINT)|(1<<TWSTO)|(1<<TWEN);
}

uint8_t I2C_write(uint8_t v) {
	TWDR = v;
	TWCR = (1<<TWINT)|(1<<TWEN);
	while(!(TWCR & (1<<TWINT)));
	uint8_t s = TWSR & 0xF8;
	return (s == 0x18 || s == 0x28);
}

void pcf8574_autodetect(void) {
	for (uint8_t a=0x20; a<=0x27; a++) {
		I2C_start(); uint8_t ok = I2C_write((a<<1)|0); I2C_stop();
		if (ok) PCF_ADDR = a;
	}
	for (uint8_t a=0x38; a<=0x3F; a++) {
		I2C_start(); uint8_t ok = I2C_write((a<<1)|0); I2C_stop();
		if (ok) PCF_ADDR = a;
	}
}

void PCF8574_write(uint8_t b) {
	I2C_start();
	I2C_write((PCF_ADDR<<1)|0);
	I2C_write(b);
	I2C_stop();
}

void twi_lcd_cmd(const unsigned char x) {
	lcd = 0x08;
	lcd &= ~(0x01);
	PCF8574_write(lcd);
	twi_lcd_4bit_send(x);
}

void twi_lcd_dwr(unsigned char x) {
	lcd |= 0x09;
	PCF8574_write(lcd);
	twi_lcd_4bit_send(x);
}

void twi_lcd_msg(const char *c) {
	while (*c) twi_lcd_dwr(*c++);
}

void twi_lcd_clear(void) {
	twi_lcd_cmd(0x01);
}

void twi_lcd_4bit_send(unsigned char x) {
	unsigned char temp = x & 0xF0;
	lcd &= 0x0F;
	lcd |= temp;
	lcd |= 0x04;
	PCF8574_write(lcd);
	_delay_us(1);
	lcd &= ~0x04;
	PCF8574_write(lcd);
	_delay_us(5);

	temp = (x & 0x0F)<<4;
	lcd &= 0x0F;
	lcd |= temp;
	lcd |= 0x04;
	PCF8574_write(lcd);
	_delay_us(1);
	lcd &= ~0x04;
	PCF8574_write(lcd);
	_delay_us(5);
}

void twi_lcd_init(void) {
	pcf8574_autodetect();
	lcd = 0x04;
	PCF8574_write(lcd);
	_delay_us(25);
	twi_lcd_cmd(0x03);
	twi_lcd_cmd(0x03);
	twi_lcd_cmd(0x03);
	twi_lcd_cmd(0x02);
	twi_lcd_cmd(0x28);
	twi_lcd_cmd(0x0F);
	twi_lcd_cmd(0x01);
	twi_lcd_cmd(0x06);
	twi_lcd_cmd(0x80);
	twi_lcd_msg("Initializing...");
	_delay_ms(1000);
	twi_lcd_clear();
	twi_lcd_cmd(0x80);
}

void lcd_twolines(const char *c, const char *b){
	twi_lcd_cmd(0x80);
	twi_lcd_msg(c);
	twi_lcd_cmd(0xC0);
	twi_lcd_msg(b);
}

\end{lstlisting}

\paragraph{Inicializaciones}\label{anexo:initsSPI}---\\
\begin{lstlisting}[language=C, caption={Inits en spi}]


void Init_timer1(void) {
	TCCR1A=0;
	TCCR1B=(1<<ICNC1)|(1<<ICES1)|(1<<CS11);
	TCNT1=0;
	TIMSK1=(1<<ICIE1);
	sei();
}

void timer2_init(void){
	TCCR2A = (1 << WGM21);
	TCCR2B = (1 << CS22)|(1<<CS21)|(1<<CS20);
	OCR2A  = 255;
	TIMSK2 = (1 << OCIE2A);
}

void Init_adc(void) {
	ADCSRA=(1<<ADEN)|(1<<ADPS2)|(1<<ADPS1)|(1<<ADPS0);
	DIDR0=(1<<ADC0D)|(1<<ADC1D);
}

uint16_t adc_read_AC0(void) {
	ADMUX=(1<<REFS0);
	ADCSRA|=(1<<ADSC);
	while(ADCSRA&(1<<ADSC));
	return ADC;
}

uint16_t adc_read_AC1(void) {
	ADMUX=(1<<REFS0)|(1<<MUX0);
	ADCSRA|=(1<<ADSC);
	while(ADCSRA&(1<<ADSC));
	return ADC;
}


\end{lstlisting}

